\chapter{Verifying Compiler Transformations}
\label{ch:verifying-compiler-transformations}

Replacing a multiplication by an addition can preserve every result while moving a division across a branch can introduce a new trap.
Both changes may leave a well-typed, well-formed control-flow graph.
A structural verifier cannot tell them apart: it checks the language of the IR, not the relationship between the old and new programs.
This chapter turns that relationship into explicit proof obligations.
It connects the analyses of Chapter~\ref{ch:designing-sound-analyses}, the solvers of Chapter~\ref{ch:automated-reasoning-limits}, and the state exploration of Chapter~\ref{ch:model-checking-algorithms} to the decision to accept an optimization.
The goal is a validator that explains both its successes and its limits, rather than a second optimizer that silently repeats the first one's mistakes.

\section{Choose the behavior before choosing the checker}
\label{sec:ch47-behavior}

Consider an eight-bit unsigned intermediate language.
Its values are $0$ through $255$, addition and multiplication wrap modulo $256$, and division by zero produces a defined trap.
The following three candidates pose different validation problems:
\begin{enumerate}
\item Replace $2x \bmod 256$ by $(x+x)\bmod 256$.
\item Replace the unsigned comparison $((x+1)\bmod 256)>x$ by true.
\item Move $18/d$ out of a conditional whose false arm returns zero.
\end{enumerate}
The first is valid for every value in the domain.
The second fails at $x=255$: the original comparison is $0>255$, hence false.
The third fails when the condition is false and $d=0$: the original returns zero and the moved operation traps.
Neither the mathematical-integer identity $x+1>x$ nor the usual algebra of rational division describes this language.
The validator must encode the selected language, including its exceptional outcomes.

For a deterministic terminating fragment, define an observation as one of
\[
  \operatorname{return}(v,\tau)
  \quad\text{or}\quad
  \operatorname{trap}(k,\tau),
\]
where $v$ is the result, $k$ is a trap kind, and $\tau$ is the ordered sequence of observable events already emitted.
A trap payload is not a returned value.
Two observations are equal only if their constructors, relevant fields, and event sequences agree.
If memory is observable, include the relevant final memory or a specified projection of it in the returned observation.
Ignoring all memory while validating a store-elimination pass would change the question being proved.

This observation model includes no undefined behavior.
A source language that has undefined behavior needs a separate account of which executions impose obligations.
For example, a validator can restrict a rewrite theorem to source-defined inputs, while still checking that the target does not introduce an error on those inputs.
It must not classify a defined language trap as undefined behavior merely because both terminate an execution abnormally.
The C draft distinguishes constraints, undefined behavior, and observable behavior; an IR validator needs its own corresponding definitions~\cite{wg14-2011-c11}.

For nondeterministic programs, a single returned value is no longer the entire specification.
Let $\mathcal{B}_S(i)$ and $\mathcal{B}_T(i)$ be the permitted source and target behaviors on input $i$ under a fixed environment model.
One common \term{refinement} obligation is
\[
  \mathcal{B}_T(i)\subseteq\mathcal{B}_S(i)
\]
on the inputs for which the source specification imposes that obligation.
Equality is a stronger requirement and may exclude legitimate reductions of nondeterminism.
Which traces count, how termination is observed, and how undefined behavior is treated must be fixed before this inclusion has a useful meaning.
This chapter's executable examples deliberately use the simpler deterministic, terminating semantics.

\section{Three different objects can be verified}
\label{sec:ch47-three-objects}

A pass implementation, a particular translation, and an abstract model of a compiler service are different verification targets.
Proving the pass implementation correct establishes a theorem about every accepted input, within the theorem's assumptions.
\term{Translation validation} instead checks the particular input/output pair produced by a run of the compiler~\cite{pnueli-1998-validation}.
Model checking a cache-invalidation protocol explores a model of the compiler's own operation; its safety result applies to that model and transfers to the implementation only through a justified modeling relation.
Table~\ref{tab:ch47-verification-targets} keeps these scopes separate.

\begin{table}[htbp]
\centering
\caption{The object being checked determines what a successful verification establishes.}
\label{tab:ch47-verification-targets}
\begin{tabular}{p{0.21\linewidth}p{0.34\linewidth}p{0.34\linewidth}}
\toprule
Target & Established obligation & Remaining connection \\
\midrule
Pass implementation & Every admitted translation preserves its contract & Implementation, semantic model, and proof assumptions agree \\
One translation & This source/target pair satisfies the encoded relation & Other translations need their own validation \\
Compiler-service model & All modeled executions satisfy the property & Real service behavior is represented by the model \\
Test executions & These executions agree with the oracle & Untested inputs and behaviors remain open \\
\bottomrule
\end{tabular}
\end{table}

The third target is the focus of Chapter~\ref{ch:model-checker-case-study}.
The second is particularly convenient at an optimizer boundary.
The compiler proposes a changed region; a validator either accepts it or retains the old region.
Rejecting a correct candidate costs an optimization opportunity.
Accepting an incorrect candidate changes program behavior.
That asymmetry justifies conservative failure when the validator cannot express an operation or a solver exhausts its resources.
Figure~\ref{fig:ch47-validation-gate} shows the handoff.

\begin{figure}[htbp]
\centering
\begin{tikzpicture}[node distance=10mm and 12mm]
  \node[basic block] (old) {original region};
  \node[basic block,right=of old] (opt) {optimizer};
  \node[basic block,right=of opt] (new) {candidate};
  \node[basic block,below=of opt] (check) {validator};
  \node[basic block,below left=of check] (keep) {keep original};
  \node[basic block,below right=of check] (accept) {install candidate};
  \draw[cfg edge] (old) -- (opt);
  \draw[cfg edge] (opt) -- (new);
  \draw[cfg edge] (old) |- (check);
  \draw[cfg edge] (new) |- (check);
  \draw[cfg edge] (check) -- node[left,font=\scriptsize] {failed or inconclusive} (keep);
  \draw[cfg edge] (check) -- node[right,font=\scriptsize] {proved} (accept);
\end{tikzpicture}
\caption{A candidate reaches installation only after its relation to the retained original has been checked.}
\label{fig:ch47-validation-gate}
\end{figure}


The optimizer can suggest a correspondence between variables, blocks, and memory locations.
Those suggestions are useful witnesses, not trusted facts.
The validator checks every premise it actually uses, including initial-state correspondence, operand availability, and exceptional exits.
Otherwise an incorrect optimizer can supply an incorrect witness that persuades an equally trusting validator.

\section{A local validator with explicit outcomes}
\label{sec:ch47-local-validator}

Start with acyclic regions containing fixed-width arithmetic, conditionals, and the defined division trap.
Each region terminates, and its operations have exact symbolic encodings.
Both regions receive the same logical inputs; temporary names are renamed apart.
Let $P(i)$ be the candidate's precondition, and let $O_S(i)$ and $O_T(i)$ encode the two observations.
The counterexample formula is
\begin{equation}
  P(i)\land\bigl(O_S(i)\neq O_T(i)\bigr).
  \label{eq:ch47-counterexample}
\end{equation}
For these total, deterministic encodings, unsatisfiability establishes equality for every input satisfying $P$.
Satisfiability provides an input at which the observations disagree.
An unknown result establishes neither statement.

The encoder represents each observation by a tag plus payloads.
Its mismatch predicate compares payloads only when the tags make them meaningful.
For example, two zero-division traps with equal event histories do not disagree because their unused result-bit fields differ.
Conversely, equal result bits do not establish agreement between a return and a trap.
An ordered event history can be encoded with a bounded event array and length for a bounded acyclic region, but the array bound must cover every possible path through that region.

Algorithm~\ref{alg:ch47-validate-local} returns a proof result distinct from a tested counterexample and a failure to decide.
Replaying a solver model with independent concrete semantics catches many encoding mistakes, although successful replay cannot establish that the encoding omitted no other behavior.

\begin{algorithm}[htbp]
\caption{Validate a deterministic, terminating local rewrite.}
\label{alg:ch47-validate-local}
\begin{pseudocode}
\Require acyclic regions $S,T$, input precondition $P$, exact operation semantics
\Ensure valid under $P$, a replayed counterexample, or inconclusive
\State check structural validity, input alignment, and supported operations
\If{a required check fails}
  \State \Return inconclusive with the failed check
\EndIf
\State construct symbolic observations $O_S(i)$ and $O_T(i)$
\State $q\gets P(i)\land(O_S(i)\neq O_T(i))$
\State $r\gets\Call{Solve}{q}$
\If{$r$ is UNSAT}
  \State \Return valid under $P$ and the encoded semantics
\ElsIf{$r$ is SAT with assignment $a$}
  \State evaluate $P(a)$, $S(a)$, and $T(a)$ using concrete semantics
  \If{$P(a)$ holds and the concrete observations differ}
    \State \Return counterexample $a$ and both observations
  \EndIf
  \State \Return inconclusive: encoding or replay disagreement
\Else
  \State \Return inconclusive with the solver status
\EndIf
\end{pseudocode}
\end{algorithm}

For $N$ source and target instructions, a shared symbolic expression DAG avoids expanding the same definition once per use.
Encoding cost must include the widths and operations, not merely count IR instructions.
If bit-blasting produces $B$ Boolean nodes and clauses, the construction and storage can be linear in $B$ with a bounded-arity DAG encoding.
The solver cost is additional and has no general polynomial bound for arbitrary Boolean inputs.
Concrete replay is linear in the executed instruction and event counts under the stated fixed-width operation model.

This algorithm is a design for a solver-backed validator, not a claim that a general SMT integration accompanies this book.
The finite-domain implementation below evaluates the same equality obligation exhaustively for its small examples.
It has no unknown solver result because it enumerates every member of an explicitly bounded domain.
Applying its conclusion to a wider integer type requires a new proof or enumeration, not a change to the report heading.

\section{Work the precondition, not just the identity}
\label{sec:ch47-worked}

For the doubling rewrite, the two formulas are
\[
  O_S(x)=\operatorname{return}((2x)\bmod256,\epsilon),\qquad
  O_T(x)=\operatorname{return}((x+x)\bmod256,\epsilon).
\]
They agree by the distributive definition of multiplication by two over integers followed by the same reduction.
At $x=137$, both return $18$, because $274=256+18$.
No no-overflow premise is necessary for this wrapping language.

For the comparison rewrite, $x=255$ gives a counterexample and every $x<255$ satisfies the proposed identity.
Thus $P(x)\equiv x<255$ is sufficient.
It is also necessary if the replacement is the constant true over the entire byte domain.
A range analysis may establish $P$ at one program point, but that proof belongs to that point and program revision.
If an earlier transformation introduces a new predecessor carrying $255$, the old range fact no longer justifies the replacement.
The validator must either derive $P$ again or verify an analysis certificate against the changed graph.

For the hoist, define $c\in\{0,1\}$ and $d\in\{0,\ldots,255\}$.
The source returns $18/d$ if $c=1$, trapping if the divisor is zero, and returns $0$ otherwise.
The target evaluates $18/d$ first and then chooses between that value and zero.
At $(c,d)=(0,0)$ their observations differ.
On $c=1$, both evaluate the division; on $d\neq0$, moving this pure, nontrapping division cannot create an event or trap in our fragment.
The exact validity region for this particular rewrite is therefore
\[
  P(c,d)\equiv(c=1)\lor(d\neq0).
\]
This is more precise than requiring $d\neq0$ unconditionally.
An implementation may choose that stronger precondition because its analysis can establish it cheaply.
Rejecting a rewrite when the sufficient condition is unknown is preferable to silently treating unknown as true.

Table~\ref{tab:ch47-finite-results} reports the finite checks performed by the accompanying C11 program.
The empty event trace is fixed in these examples; a pass that moves effects requires a richer observation comparator.

\begin{table}[htbp]
\centering
\caption{Exhaustive checks distinguish unconditional identities from guarded rewrite permissions.}
\label{tab:ch47-finite-results}
\begin{tabular}{p{0.42\linewidth}rrp{0.24\linewidth}}
\toprule
Rewrite & Inputs & Mismatches & First counterexample \\
\midrule
Wrapping doubling to addition & 256 & 0 & None \\
Unsigned successor comparison to true & 256 & 1 & $x=255$ \\
Hoist conditional division & 512 & 1 & $c=0,d=0$ \\
Hoist under $(c=1)\lor(d\neq0)$ & 511 & 0 & None \\
\bottomrule
\end{tabular}
\end{table}

The implementation in \path{example-code/translation-checks/} keeps source and target evaluation functions separate and compares tagged outcomes.
It enumerates both the candidate's full domain and the guarded subset, asserting the exact mismatch counts.
It checks the derivation above rather than supplying evidence about a production compiler or an arbitrary transformation language.
Chapter~\ref{ch:testing-validating-compilers} explains how to retain such witnesses as regression tests.

\section{Loops require a relation between executions}
\label{sec:ch47-simulation}

Instruction-by-instruction equality fails for transformations that change control structure.
Unrolling one loop iteration into two source iterations changes the number of loop-header visits.
A \term{simulation relation} can instead associate states at selected semantic boundaries.
It relates values, memory, and control positions that represent the same progress, rather than demanding equal program counters.
Figure~\ref{fig:ch47-cutpoints} illustrates the grouping.

\begin{figure}[htbp]
\centering
\begin{tikzpicture}[x=23mm,y=13mm]
  \node[basic block] (s0) at (0,1.5) {source $i$};
  \node[basic block] (s1) at (2,1.5) {$i+1$};
  \node[basic block] (s2) at (4,1.5) {$i+2$};
  \node[basic block] (t0) at (0,0) {target $j$};
  \node[basic block] (t2) at (4,0) {$j+2$};
  \draw[cfg edge] (s0) -- node[above,font=\small] {one iteration} (s1);
  \draw[cfg edge] (s1) -- node[above,font=\small] {one iteration} (s2);
  \draw[cfg edge] (t0) -- node[below,font=\small] {one unrolled iteration} (t2);
  \draw[dashed] (s0) -- node[left,font=\small] {$R$} (t0);
  \draw[dashed] (s2) -- node[right,font=\small] {$R$} (t2);
\end{tikzpicture}
\caption{A simulation relates paired boundaries; two source iterations correspond to one target iteration without equating intermediate program counters.}
\label{fig:ch47-cutpoints}
\end{figure}


For a concrete example, let the source sum $0+1+\cdots+(n-1)$ with mathematical integers, initial $i=s=0$, and fixed nonnegative $n$.
Its step executes $s\gets s+i$ followed by $i\gets i+1$ whenever $i<n$.
The target uses an unrolled step only when $j+1<n$:
\[
  t\gets t+j+(j+1),\qquad j\gets j+2.
\]
After that loop, a scalar cleanup step runs if $j<n$.
At paired cutpoints use $R(i,s,j,t)\equiv i=j\land s=t$.
One unrolled target step changes $t$ to $t+2j+1$ and $j$ to $j+2$.
Two source steps change $s$ to $s+i+(i+1)$ and $i$ to $i+2$.
Substituting $i=j,s=t$ proves $R$ at the next paired point.
The guard $j+1<n$ ensures both source steps are enabled.
The cleanup step matches one source step, and both exits imply $i=j=n$ and $s=t$.

For $n=5$, the paired states $(i,s)=(j,t)$ are $(0,0)$, $(2,1)$, $(4,6)$, then $(5,10)$ after cleanup.
The source also visits $(1,0)$ and $(3,3)$ between the paired states.
Those intermediate control points are deliberately absent from the relation; matching them to arbitrary target headers would create false obligations.
The C lab checks both executions at every actual paired boundary for $0\leq n\leq64$, alongside the mathematical proof just given.
That finite test range is a bug check, not the reason the unbounded mathematical theorem holds.

Stuttering, where one side takes internal steps while the other stays at a boundary, needs a progress argument if termination or infinite observations matter.
Without one, a target that loops forever taking unobservable steps might appear to match a source that returns immediately.
Here $n-j$ is a nonnegative ranking expression at each enabled target loop step and decreases by two, or by one in cleanup.
The source ranking expression $n-i$ decreases by one.
Both terminate for finite nonnegative $n$, so equality at exits yields equality of their results.
For fixed-width counters the guard and arithmetic require additional overflow checks; this mathematical proof alone cannot justify a machine-word loop rewrite.

\section{Check analysis certificates independently}
\label{sec:ch47-certificates}

An optimizer often already has facts that could make validation cheaper.
It may know a loop invariant, a range bound, or that a pointer targets one object.
Treat these facts as a candidate certificate and verify their semantic obligations.
Computing an invariant can be expensive and heuristic; checking a proposed invariant can be a finite set of implication queries.
This separation resembles the distinction between discovering a proof and checking one.

Let $I_B$ be a state predicate proposed at entry to block $B$, and let $T_e(\sigma,\sigma')$ encode the semantics of an edge including its block execution and guard.
An entry obligation requires the initial-state predicate to imply $I_{\Entry}$.
Each edge requires
\[
  I_B(\sigma)\land T_{B\to C}(\sigma,\sigma')
       \Longrightarrow I_C(\sigma').
\]
A use of the certificate additionally requires $I_B$ to imply the specific safety or rewrite precondition at the point where it is used.
If that point is inside $B$, compose the transfer of the preceding instructions before checking the precondition.
An invariant at block entry does not automatically survive every assignment in the block.

Algorithm~\ref{alg:ch47-check-certificate} checks the entry and consecution obligations for a fixed graph.
All implication queries are discharged by asking whether the antecedent and the negation of the consequent are satisfiable.

\begin{algorithm}[htbp]
\caption{Check a proposed invariant map before trusting its clients.}
\label{alg:ch47-check-certificate}
\begin{pseudocode}
\Require CFG with exact edge relations $T_e$, initial predicate $Init$, proposed predicates $I_B$
\Ensure certificate accepted only if initialization and every edge preserve it
\State require \Call{Solve}{$Init\land\neg I_{\Entry}$} to return UNSAT
\ForAll{edges $e=(B,C)$}
  \State $q\gets I_B(\sigma)\land T_e(\sigma,\sigma')\land\neg I_C(\sigma')$
  \State require \Call{Solve}{$q$} to return UNSAT
\EndFor
\State on any SAT, UNKNOWN, unsupported encoding, or resource failure, reject the certificate
\State \Return accepted for this CFG and these semantic parameters
\end{pseudocode}
\end{algorithm}

Induction on finite path length proves that every reachable entry state satisfies its block predicate once all obligations hold.
This is a safety invariant, not a termination theorem.
The map containing false at every block would pass many edge implications vacuously, but the initialization obligation rejects it whenever initial states exist.
The map containing true everywhere is inductive yet may imply no useful client precondition.
Both checks, initialization and client implication, are essential.

For $E$ CFG edges this validator asks $E+1$ consecution and initialization queries, followed by the client's implication queries.
The formula-construction and solver costs depend on the predicates and transfer encodings; counting queries alone is not a running-time bound.
Cache only checked obligations keyed by IR revision, semantic configuration, certificate, and relevant environment assumptions.
Changing the graph, arithmetic width, alias model, or call summary changes what was proved.

\section{A dataflow query as a model-checking problem}
\label{sec:ch47-dataflow-query}

Some compiler analyses can themselves be checked against a small state-space formulation.
Fix one variable $v$ and view CFG blocks as abstract states.
Label a block $U$ when it uses the incoming value of $v$ before redefining it, and label it $T$ when it contains no definition of $v$.
The incoming value is live at $B$ exactly when a finite CFG path reaches a $U$ block and every preceding block on that path is $T$.
For this abstract graph, that is the CTL property $E[\mathit{transparent}_v\mathbin{\mathrm{U}}\mathit{use}_v]$, using existential strong until and spelling out the two atomic labels.
The corresponding least-fixed-point equation is
\[
  L=\mu X.\bigl(U\cup(T\cap\operatorname{Pre}_{\exists}(X))\bigr).
\]
Here $U$ and $T$ also denote their sets of labeled blocks, and $\mu$ selects the least fixed point.
This is the usual backward may-liveness equation for one variable, expressed as a reachability property with a path restriction.

Take edges $e\to u$, $e\to k$, $u\to r$, and $k\to r$.
Block $u$ reads $v$, block $k$ assigns a fresh value to $v$, and block $r$ returns $v$.
Block $e$ branches without reading or defining $v$.
Then $U=\{u,r\}$ and $T=\{e,u,r\}$.
Starting with the empty set gives successive approximations $\emptyset$, $\{u,r\}$, and $\{e,u,r\}$, after which nothing changes.
The old value is not live at $k$: the read at $r$ sees the new definition.
It is live at $e$ because the path through $u$ uses it before any redefinition.

This formulation gives an independent oracle for a bit-vector liveness implementation on small graphs.
It is exact for the CFG path abstraction, not necessarily for concrete executions: an infeasible branch can still supply a CFG path.
Consequently, a positive answer is conservative evidence of possible use, not a concrete input that executes that use.
A negative answer permits discarding the old value only with respect to the modeled uses; exceptional edges, deoptimization recovery, and other observable consumers must also be represented.
This example connects dataflow and model checking through their shared fixed-point equations without claiming that every dataflow framework is reducible to this particular CTL formula~\cite{nielson-1999-principles,baier-2008-model-checking}.
Schmidt develops the broader relationship using abstract interpretations and modal fixed-point properties~\cite{schmidt-1998-dataflow-model-checking}.

\section{Implementation boundaries that affect trust}
\label{sec:ch47-trust}

The validator still depends on a trusted semantic encoder and a solver or proof checker.
An UNSAT certificate can reduce reliance on the solver's search implementation when a separate checker validates it, but only for the supported proof language and theories.
It does not establish that the encoder correctly modeled the input IR.
Cross-check encoding rules with an independent interpreter on small domains, including division, shifts, exceptions, and conditional execution.
Keep those tests separate from optimization tests so a shared mistake is easier to locate.

Memory adds distinctions absent from scalar identities.
Two loads from the same symbolic address need not return the same value if an intervening store can alias it.
A call may modify reachable memory even when its result is unused.
Removing a volatile read can change the event trace while leaving the returned integer unchanged.
Do not import an unproved alias assertion from the optimizer into the validation formula; establish it from the modeled memory or verify its certificate.
Chapter~\ref{ch:interprocedural-heap-analysis} supplies the abstractions on which such certificates can be based.

Concurrency requires an explicit memory model and matching observations, not merely exploring all sequential interleavings of source instructions.
An allowed hardware behavior under a relaxed memory model need not appear in that sequentially consistent model.
The finite-state techniques of Chapter~\ref{ch:transition-temporal-properties} still apply to a chosen finite transition system, but the transition relation must encode the intended semantics.
Unsupported concurrency or memory operations should produce an explicit unsupported result, not silently weaken the property.

Production use also needs a failure policy.
Retain an unmodified region or a valid baseline artifact until its replacement passes validation.
A timeout records the candidate, precondition, semantic version, and solver status; it need not stop compilation if the original code is available.
A replayed counterexample records the input and both observations and becomes a regression case.
These policies keep validation failures diagnosable while allowing the optimizer's search to remain heuristic.

\section{Exercises}
\label{sec:ch47-exercises}

\begin{enumerate}
\item For the byte-domain rewrite $((x+1)\bmod256)>x\rightsquigarrow\operatorname{true}$, give the exact set of valid inputs, the sole counterexample, and a sufficient interval fact at the use site.\label{ex:ch47-byte-guard}
\item Change the conditional-division example to emit an event before the division in the true arm, and hoist both operations before the branch. Describe a counterexample with nonzero $d$ that a returned-value-only comparator misses.\label{ex:ch47-event-hoist}
\item Repeat the paired-loop trace for $n=6$. List all paired states and state why the cleanup does not execute. Explain which bound in the proof would fail if unsigned subtraction wrapped before the loop test.\label{ex:ch47-loop-pairs}
\item A certificate has $I_B\equiv\operatorname{true}$ for all blocks. Which obligations of Algorithm~\ref{alg:ch47-check-certificate} can it satisfy, and why does this not justify division-check elimination?\label{ex:ch47-vacuous-certificate}
\end{enumerate}
Answers and checks are in \path{example-code/translation-checks/}.

\section{Further reading}
\label{sec:ch47-reading}

Pnueli, Siegel, and Singerman formulate translation validation as checking the result of a translation rather than proving the translating implementation~\cite{pnueli-1998-validation}.
Baier and Katoen develop simulation, trace properties, and abstraction, which supply the vocabulary needed to state exactly what a state-space check preserves~\cite{baier-2008-model-checking}.
Nielson, Nielson, and Hankin connect analysis frameworks to their semantic justification; use that distinction when deciding which inferred facts are safe to turn into rewrite preconditions~\cite{nielson-1999-principles}.
Huth and Ryan explain logical validity, program reasoning, and model checking with explicit proof obligations~\cite{huth-2004-logic}.
